\section{Closed form of the six clinical properties}
\label{sec:supp_predicates}
Table~\ref{tab:predicates_full} gives the score and the failure map of every property, and the
expanded definitions of every subterm follow. All quantities are computed from the noisy input and
the backbone output alone. No clean reference is used at any point.

\begin{table*}[t]
\centering
\caption{Closed form of the six clinical properties. For each predicate, the table lists its clinical role, image-level score, and pixel-level failure map. larger $\mathbf{m}_i$ means stronger evidence that the corresponding property is locally unsatisfied. The expanded subterm definitions follow immediately below. Here $\bar e=\operatorname{mean}(e(x))$, $b_e=\mathbf{1}[e(x)>\bar e]$, $v(z)=\operatorname{mean}_w z$ is a depth profile, $\rho_h$ is 1-pixel horizontal autocorrelation, and $\operatorname{Var}_k$ is local variance over a $k\times k$ window.}
\label{tab:predicates_full}
\setlength{\tabcolsep}{3pt}
\footnotesize
\begin{tabular}{cp{0.18\textwidth}p{0.33\textwidth}p{0.28\textwidth}c}
\toprule
$P_i$ & Clinical role & Aggregate score $s_i$ & Failure map $\mathbf{m}_i$ & $\tau_i$ \\
\midrule
$P_1$ & Edge continuity &
$0.4\,c_{\text{cont}}+0.3\,c_{\text{efr}}+0.3\,c_{\text{ang}}$ &
$(1-\tilde e)(1-b_e)$, \ $\tilde e=e(x)/(e_{\max}+\varepsilon)$ &
0.6 \\
$P_2$ & Local contrast recovery &
$0.20\,s_{\text{range}}+0.35\,s_{\text{cnr}}+0.20\,s_{\text{var}}+0.25\,s_{\text{cov}}$ &
$[\operatorname{ReLU}(\bar\sigma-\sigma_9(x))/(\bar\sigma+\varepsilon)+0.3\,\mathbf{1}[\sigma_9(x)<0.01]]_{[0,1]}$ &
0.5 \\
$P_3$ & Flat-region smoothness &
$0.6\,q_{\text{var}}+0.4\,q_{\text{res}}$ &
$F \odot \sigma_5(x)/(\max \sigma_5(x)+\varepsilon)$ &
0.75 \\
$P_4$ & Structural coherence &
$0.4\,s_{\text{self}}+0.3\,s_{\text{reg}}+0.3\,s_{\text{vert}}$ &
$e(e(x))/(\max e(e(x))+\varepsilon)$ &
0.5 \\
$P_5$ & Speckle conformity &
$0.7\,\operatorname{mean}(\exp(-|cv-cv^\star|/0.15))+0.3\,\exp(-5|\rho_h(z)|)$ &
$1-\exp(-|cv-cv^\star|/0.15)$ &
0.3 \\
$P_6$ & Layer visibility &
$0.4\,s_{\text{clr}}+0.3\,s_{\text{dep}}+0.3\,s_{\text{smo}}$ &
$\operatorname{Var}_7(e(x))/(\max \operatorname{Var}_7(e(x))+\varepsilon)$ &
0.7 \\
\bottomrule
\end{tabular}
\end{table*}

{\footnotesize
\noindent\textbf{Expanded subterms in Table~\ref{tab:predicates}.}
\textbf{$P_1$.}
$c_{\text{cont}}=\langle\operatorname{erode}(\operatorname{dilate}(b_e)),b_e\rangle/(\|b_e\|_1+\varepsilon)$,
$c_{\text{efr}}=\exp(-5|r_e-0.2|)$ with
$r_e=\|\mathbf{1}[e(x)>1.5\bar e]\|_1/(\|\mathbf{1}[e(x)<0.5\bar e]\|_1+\varepsilon)$,
and $c_{\text{ang}}=\exp(-2\,\operatorname{mean}(\sigma_5(\theta(x))))$.
\textbf{$P_2$.}
$s_{\text{range}}=\operatorname{sigm}(10(\Delta_r-0.3))$, $\Delta_r=\max(x)-\min(x)$.
$s_{\text{cnr}}=\operatorname{sigm}(15(\bar\sigma-0.05))$, $\bar\sigma=\operatorname{mean}(\sigma_9(x))$.
$s_{\text{var}}=\max(0.2,\operatorname{sigm}(-20(v_\sigma-0.02)+2))$, $v_\sigma=\operatorname{Var}(\sigma_9(x))$.
$s_{\text{cov}}=\operatorname{sigm}(5(\gamma-0.4))$, $\gamma=\operatorname{mean}(\mathbf{1}[\sigma_9(x)>0.3\bar\sigma])$.
\textbf{$P_3$.}
$F=\mathbf{1}[e(x)<0.3\bar e]$,
$q_{\text{var}}=1-\langle F,\sigma_5(x)/(\sigma_5(y)+\varepsilon)\rangle/(\|F\|_1+\varepsilon)$,
and $q_{\text{res}}=\exp(-5\,\operatorname{mean}|y-x|)$.
\textbf{$P_4$.}
$s_{\text{self}}$ is the average Pearson correlation between $x$ and shifts $(0,2)$, $(2,0)$, and $(2,2)$.
$s_{\text{reg}}=\exp(-12\,\operatorname{mean}(e(e(x))))$.
$s_{\text{vert}}=\operatorname{clip}(\max |\nabla_z v(x)|/(8\,\operatorname{mean}|\nabla_z v(x)|),0,1)$.
\textbf{$P_5$.}
$z=\log(y+\varepsilon)-\log(x+\varepsilon)$,
$cv=\sigma_{15}(|z|)/(\mu_{15}(|z|)+\varepsilon)$,
and $cv^\star(d)$, where $d$ is the axial depth index, is the depth-wise Gamma-K target coefficient of variation used by the implementation.
\textbf{$P_6$.}
$s_{\text{clr}}=\operatorname{mean}(\operatorname{clip}(\max |\nabla_z v(x)|/(10\,\operatorname{mean}|\nabla_z v(x)|),0,1))$,
$s_{\text{dep}}=\operatorname{clip}(20\,\operatorname{Var}(v(x)),0,1)$,
and $s_{\text{smo}}=1-\operatorname{clip}(10\,\operatorname{mean}(\operatorname{Var}_z(\operatorname{mean}_w e(x))),0,1)$.

These are the implemented predicates. The maps $\mathbf{m}_i$ are the spatial quantities consumed by the negotiator and potential head, while the scores $s_i$ drive the rule activations.
}
\noindent The released code computes exactly these quantities. The relevant class is named
\texttt{EnhancedGTFreePredicates} and the speckle property is computed by
\texttt{PhysicsAccurateSpecklePredicate}.
